\documentclass[11pt]{article}

% "final" = camera-ready (author block shown, no line numbers).
% "review" = anonymised submission with line numbers.
% "preprint" = non-anonymous version with page numbers.
\usepackage[final]{acl}
\usepackage{times}
\usepackage{latexsym}
\usepackage[T1]{fontenc}
\usepackage[utf8]{inputenc}
\usepackage{microtype}
\usepackage{inconsolata}
\usepackage{booktabs}
\usepackage{amsmath}
\usepackage{amssymb}
\usepackage{graphicx}

% ============================================================
% TITLE AND AUTHORS
% ============================================================

\title{EcoExplain-VLM: Evaluating Vision--Language Models for Classifying
and Explaining Forest Vegetation Change in Kanha Tiger Reserve}

\author{
  Rama Chandra Prasad Pillutla \and
  Chinta Phani Rama Vaibhav \and
  Shaik Afzal Basha \\
  Lab for Spatial Informatics, International Institute of Information Technology Hyderabad, India \\
  \texttt{rcprasad@iiit.ac.in}, \texttt{chinta.vaibhav@students.iiit.ac.in}, \\
  \texttt{afzalbasha.shaik@students.iiit.ac.in}
}

\begin{document}

\maketitle

% ============================================================
% ABSTRACT
% ============================================================

\begin{abstract}
Observation of changes in vegetation in protected forests requires an approach which combines information obtained from a variety of sources. This paper shows how small and reproducible vision-language setup can be used to introduce this evidence into open vision-language models (VLMs) and evaluate their performance with respect to both classification results and soundness of their reasoning. The dataset for this purpose is called EcoExplain-VLM and it consists of 366 forest patches in Kanha Tiger Reserve in Madhya Pradesh, India. Pairs of before/after images taken during dry season in January-February 2022 and 2023 from Sentinel-2, Sentinel-1 SAR, CHIRPS, ERA5-Land, and NASADEM are presented along with corresponding evidence card.Training, development, and test data are distinguished (256/55/55). On the same dataset of 55 test patches, the performance of Qwen2-VL-2B in zero-shot setting is 56.4\%, which is 21.9\% higher than the majority class baseline of 34.5\%. In contrast, the classification score of LLaVA-Interleave-0.5B, PaliGemma-3B, and BLIP-2 OPT-2.7B lies in the range of 32.7-36.4\%. A single threshold of change in NDVI, estimated on training data, leads to 94.5\% correct classifications, thus indicating the high level of recoverability of the spectral rank and the difficulty of the task for language models. The use of Low-Rank Adaptation (LoRA) with 16 training patches and one epoch of fine-tuning of Qwen2-VL-2B improves its accuracy up to 67.3\%. The accuracy scores of Qwen2-VL-2B are given for the entire evidence card, with the last summary line containing the candidate label. If that last line is excluded from the card, the models can be steered to produce evidence-based text: the cite-only prompt brings claim correctness to 75.8\% for Qwen2-VL-2B and 71.1\% for LLaVA-Interleave-0.5B, while free form prompts lead to little verifiable claims or those contrary to the measurements.A pre-trained LoRA for sixteen self-generated explanations filtered achieves 61.8\% class accuracy in the explanation mode (where 50.5\% are accurate claims), which means that the explanation path can be trained on this split of data. EcoExplain-VLM provides thus a productive way forward: a leakage-aware, multi-sensor testbed appropriate for model comparison at scale or for adaptation, and a prompting strategy that keeps alignment of open VLMs with the provided facts.
\end{abstract}

\noindent
\textbf{Keywords:}
vision-language models;
forest vegetation change;
multi-source remote sensing;
Sentinel-2;
Sentinel-1 SAR;
evidence card;
zero-shot classification;
parameter-efficient fine-tuning;
LoRA;
evidence-grounded explanation;
Kanha Tiger Reserve

% ============================================================
% 1. INTRODUCTION
% ============================================================

\section{Introduction}

State Forestry Departments, which operate under the governance of their own states in collaboration with national level monitoring organizations,
should be aware not only whether the greenness or browning of canopy occurred between
two time periods, but also what made it reliable in these reserved areas. Optical indices like NDVI, backscatter from radar, rainfall, soil moisture anomalies, and topography are all contributing together in this regard.
Forest monitoring is facilitated by large-scale operation change products, where it is possible to enhance their results through adding natural language explanations allowing verification of interpretation in relation to the original data.

Vision-language models (VLMs) map both images and text to a common representation space, and produce unstructured text outputs. Such a model in principle can analyze two chips of dry season and an evidence card in a concise format, and give one of the following labels
(\texttt{stable}, \texttt{decline}, \texttt{increase}) together with a
justification. In reality, small open models hallucinate content, produce explanations that dont reflect how they
reached their answer and are sensitive to
prompt phrasing. Remote-sensing VLMs such as GeoChat are trained for overhead imagery and spatial referring expressions. The setting here is different: matched dry-season chips plus a multi-sensor forest-change card, with both a class decision and a checkable justification.

This study asks three questions for a single, compact reserve:
\begin{enumerate}
\item Can open VLMs with 0.5--3 billion parameters classify relative
vegetation change from before/after RGB chips and an evidence card
without task-specific training?
\item How does a small LoRA update change that
behaviour?
\item With the label removed from the card, do the models' written
justifications agree with the measured changes, and does prompt
structure or fine-tuning on justifications improve that agreement?
\end{enumerate}

The contribution is a fully specified benchmark, EcoExplain-VLM, for
Kanha Tiger Reserve, together with an evaluation that reports its own
weaknesses. On the 55-patch test set, only Qwen2-VL-2B exceeds the
majority-class baseline (56.4\% against 34.5\%), and a 16-patch LoRA
raises it to 67.3\%. Both figures use an evidence card that ends with a
pipeline summary line containing the candidate label. The other three
models answer almost every patch with the same class, while a
one-number NDVI-change rule reaches 94.5\%. For explanations, a
structured ``cite-only'' prompt raises claim correctness to the 70\%
range, but the resulting texts restate the supplied facts rather than
interpret the images. Fine-tuning Qwen2-VL-2B on 16 of its own filtered
justifications gives 61.8\% class accuracy and 50.5\% correct claims.

% ============================================================
% 2. DATA AND METHODS
% ============================================================

\section{Data and Methods}

\subsection{Study Area}

Kanha Tiger Reserve lies in the Maikal hills of central India, in the
Mandla and Balaghat districts of Madhya Pradesh.The core area spans an area of
917.44\,km$^2$ whereas the buffer area is 1,134.36\,km$^2$, making a total
of 2,051.79\,km$^2$. The forest is primarily made up of sal (\emph{Shorea
robusta}) and mixed deciduous forests with grassland meadows. Rain mostly
falls in the period of June-September whereas it is cool and dry in
January-February.

Processing made use of the rectangle
$80.433^\circ$--$81.050^\circ$\,E,
$22.117^\circ$--$22.450^\circ$\,N (2,350.7\,km$^2$ in Earth Engine).
The rectangle is larger than the gazetted reserve, so it is possible
that some patches are situated either in the buffer zone or slightly
beyond the official boundaries. All the 1,200 candidate patches are
located within the rectangle. The pool of candidates contains an empty
stripe around $80.75^\circ$--$80.85^\circ$\,E; this happened due to
the previous repository that generated the pool.

\subsection{Data Sources}

\subsubsection{Optical and Synthetic Aperture Radar Data}

Sentinel-2 L2A surface reflectance products provide the optical bands at 10-20 m resolution. Clouds and shadows were masked using the Scene Classification Layer, and median composites were generated for the period of 1 January–28 February 2022 (T0, 20 scenes) and the same period of 2023 (T1, 23 scenes). We derived the NDVI, EVI, NDMI, and NBR from these composites and rendered true colour RGB chips for the VLMs.

Sentinel-1 GRD C-Band Backscatter products were composited for the same periods using the direction of the majority of orbit passes (descending; 15 scenes for T0 and 20 scenes for T1). The VV and VH polarizations (in dB) are sensitive to canopy structure and moisture but not affected by clouds. All exports and composites were created in Google Earth Engine.

\subsubsection{Ancillary Environmental Data}

Daily CHIRPS rainfall is added up for each of the 58-day periods.
ERA5-Land has 2\,m air temperature and volumetric soil moisture in the
top layer (0-7\,cm), provided as anomalies relative to a 2015-2021
climatology for the same months. Elevation, slope, and aspect are derived from NASADEM at 30\,m. None
of these are in the automatic label. But the numbers are needed in
the evidence card in case there is justification based on drought or
terrain.

\subsection{Preprocessing and Patch Generation}

Patch generation involves drawing a circle of 320 m radius
around each candidate center (which is half of $64 \times 10$\,m),
creating a bounding square and clipping the square against the
study rectangle, so that patches at the edges are a little bit smaller.
GeoTIFF exports are usually $70 \times 65$ pixels in EPSG:4326 (pixel size $0.0000898^\circ$).
Fifteen rasters are saved per patch for a total of 5,490 files. After cloud masking, the fraction of cloud pixels is 0.8--2.3\%
per patch; candidates with a masked fraction higher than 0.15
pixels in either year were marked as uncertain and not sampled. As
for the VLMs, the T0 and T1 RGB chips were downsampled to
$512 \times 512$ pixels.

The minimum distance between the 366 chosen centers is 639 m
(adjacent non-overlapping squares), while the median is 1,430 m.
The data split is a random stratified sampling procedure instead of a
spatial block split, so spatial autocorrelation between training
and test patches is still an issue.

\subsection{Change Evidence and Labelling}

Let $\Delta$ be the patch-mean NDVI change,
$\mathrm{NDVI}_{T1} - \mathrm{NDVI}_{T0}$. Over the 1,200-candidate
pool, the 15th and 85th percentiles of $\Delta$ are $q_{15} = -0.0377$
and $q_{85} = +0.0031$. The automatic label is
\begin{equation}
\mathrm{class} =
\begin{cases}
\text{uncertain} & \text{mask} > 0.15 \text{ in T0 or T1},\\
\text{decline} & \Delta \le q_{15},\\
\text{increase} & \Delta \ge q_{85},\\
\text{stable} & |\Delta| < 0.05,\\
\text{uncertain} & \text{otherwise}.
\end{cases}
\end{equation}
Since both percentiles fall within $\pm 0.05$, all unmasked patches from
$q_{15}$ to $q_{85}$ are classified as stable and the ``uncertain''
category will not apply to this reserve. The entire area has lost some
greenness (mean $\Delta = -0.016$, s.d.\ $0.027$) and hence ``increase''
refers to those patches that have become greener (relative) rather than
a real improvement in greenness; increase patches have $\Delta$ values of
$-0.001$ to $+0.075$ and half of them have changed by less than 0.011.
``Decline'' refers to those patches that have shown the greatest loss in
greenness ($\Delta$ from $-0.034$ to $-0.071$). Only 34 out of 366 patches
have $|\Delta| \ge 0.05$. The categories are relative rankings in this case.

The first 300 patches (100 per class) were drawn from the candidate
pool. A second batch added up to 25 patches per class, keeping new
centres at least 960\,m (1.5 patch widths) from existing ones; only 22
decline candidates passed this filter, giving 372 patches. Six patches
with incomplete exports were removed, leaving 366: 120 decline, 124
stable and 122 increase. The percentiles for the second batch were
recomputed on the remaining pool, so the class ranges overlap slightly
(decline patches reach $\Delta = -0.034$ and stable patches go down to
$-0.036$).

Every single patch comes with an evidence card with information like
indices, changes, radar measures, climate abnormalities, topography and
pipeline description. This pipeline description contains a string of
text \texttt{candidate\_class=}\ldots which is the field used for
labelling. These trend words (e.g.
\texttt{moderate\_decline}) on the evidence card have been obtained
from $512 \times 512$ tensors using thresholds and not from percentile
rule; these trend words contradict the label in 108 out of 366 cases.
The values for NDVI on the evidence card average to 0.92 of Earth
Engine mean values. Otsu thresholding bounding boxes are present in
only three patches and all narrative explanations are just templates,
because no language model based explanation was obtained during the
dataset construction.

A human reference class was assigned to 19 patches (13 training, 3
development, 3 test). It agrees with the automatic label on 8 of 19
(42\%; Wilson 95\% interval 23--64\%). An earlier
36-patch pilot agreed on 17 of 36 (47\%). There was no second
annotator. The automatic labels are therefore a reproducible ranking,
not a validated disturbance map.

\subsection{Dataset Construction and Data Splits}

Each patch is stored as one JSONL record containing image paths,
per-source statistics, the evidence card and the labels. The patches
were split 256/55/55 into training, development and test sets by
stratified sampling on the label (seed 42). The test set contains 18
decline, 18 stable and 19 increase patches, so always predicting
``increase'' gives a majority-class baseline of $19/55 = 34.5\%$. Every
result in this paper uses this test set.

\subsection{Vision--Language Model Evaluation}

Four instruction-tuned open models were compared:
Qwen2-VL-2B-Instruct ,
LLaVA-Interleave-0.5B ,
PaliGemma-3B-mix-224  and
BLIP-2 OPT-2.7B , all loaded from
Hugging Face. Parameter counts are those of the released checkpoints.
Decoding is greedy. A deterministic parser maps the output to a class:
it reads a single-word answer directly and otherwise searches the text
for the first class word or synonym (for example, ``decrease'' is read
as decline). All 55 outputs were parsed for every model.

\subsubsection{Zero-Shot Evaluation}

The main classification setting, Condition~C, gives a model both RGB
chips and the evidence card and asks for one word: \texttt{decline},
\texttt{increase} or \texttt{stable}. The four models did not receive
exactly the same input (Table~\ref{tab:inputs}). Qwen2-VL-2B and
LLaVA-Interleave-0.5B got images together with a full card, which includes
the \texttt{candidate\_class} summary line. PaliGemma would not respond
to the long prompt for analysis, hence it got images, NDVI summary line,
and a short prompt. The BLIP-2 model takes one image at a time, and hence
the two chips were put together in an image with a short prompt. Zero-shot
tests were run on a laptop.

\begin{table}[t]
\centering
\small
\begin{tabular}{lll}
\toprule
Model & Images & Text \\
\midrule
Qwen2-VL-2B & T0, T1 & full card \\
LLaVA-Interleave-0.5B & T0, T1 & full card \\
PaliGemma-3B & T0, T1 & NDVI line only \\
BLIP-2 OPT-2.7B & stitched & none \\
\bottomrule
\end{tabular}
\caption{Input per model in Condition~C. ``Full card'' includes the
pipeline summary line containing \texttt{candidate\_class}.}
\label{tab:inputs}
\end{table}

Three supporting analyses help interpret the results. First, an
\texttt{evidence\_only} variant removes the summary line and any other
line that names the label. This changes the input only for LLaVA,
because the PaliGemma and BLIP-2 prompts never contained the line, and
Qwen2-VL-2B was run in this mode on only three patches. Second, a
closed-set probe scores LLaVA by the likelihood it assigns to each class
word, with a stated decision rule and a class bias fitted on the
development set. Third, two baselines that use no model apply a
threshold to the $\Delta$NDVI value on the card: one with fixed cut-offs
of $\pm 0.02$, and one with cut-offs fitted on the training set.

\subsubsection{Parameter-Efficient Fine-Tuning}

LoRA  was applied to Qwen2-VL-2B with rank 8,
$\alpha = 16$ and dropout 0.1 on the attention and MLP projection
layers, using AdamW  with a learning rate of
$5 \times 10^{-5}$, weight decay 0.05 and gradient accumulation over
four steps. Training ran on a GTX 1080 Ti node of the institute's Ada
cluster. Two separate runs are reported.

\paragraph{Class LoRA.} The adapter was trained on 16 training patches
for one epoch with the classification target and the full Condition~C
card (final training loss 1.84), and tested on the full test set. This
is a small-scale configuration, not a complete fine-tune of the 256
training patches. Exploratory runs on all 256 training patches for two
to three epochs reached 100\% development and test accuracy, because
the model learned to copy \texttt{candidate\_class} from the summary
line. These runs are excluded from the results.

\paragraph{Explanation LoRA.} Qwen2-VL-2B itself acted as teacher: it
generated a justification for each training patch, and a generation was
kept only if its class matched the label, it contained at least two
checkable claims, and all of those claims agreed with the measured
changes. Of 256 training patches, 28 passed. A larger 7B teacher could
not be used because its download exceeded the cluster's disk quota. The
training run was capped at the first 128 training patches, so 16 of the
28 filtered justifications were used. Both training and evaluation used
the card without the label, the checkpoint was selected on the
development set, and the test set was evaluated in explanation mode.

\subsubsection{Evaluation Metrics}

Classification is scored by accuracy with Wilson 95\% intervals, macro-F1, and comparison with the 34.5\%
majority-class baseline; McNemar's test  compares
paired predictions.

Explanations are scored automatically. The scorer searches the text for
an indicator (NDVI, EVI, NBR, NDMI, rainfall, temperature, soil
moisture, VV, VH) and pairs it with the nearest direction word
(increase-type, decrease-type or stable-type). Each measured change is
converted to a direction using an indicator-specific dead-band, below
which a change counts as stable. A claim is \emph{correct} when its
direction matches the measured direction. We report the mean number of
checkable claims per explanation, the share of correct claims (pooled
over patches), the share of explanations with at least one incorrect
claim, and self-consistency (whether the dominant claimed direction
agrees with the model's own class). ROUGE-L  is
computed but not used as a measure of faithfulness, because overlap with
the supplied text rewards copying rather than reasoning. The scorer measures consistency with the evidence,
not faithfulness to the model's internal reasoning.

\subsection{Experimental Workflow}

Figure~\ref{fig:workflow} summarises the pipeline.

\begin{figure}[t]
\centering
\includegraphics[width=\columnwidth]{fig_workflow.png}
\caption{EcoExplain-VLM workflow. After the dataset is frozen, three
experiments are run on the same 55 test patches and are scored
separately.}
\label{fig:workflow}
\end{figure}

The work proceeded in six stages.

\begin{enumerate}
\item \textbf{Study design.} The study area is Kanha Tiger Reserve.
Matched dry-season windows are T0 (January--February 2022) and T1
(January--February 2023).

\item \textbf{Collect satellite and environmental data.} Sentinel-2
(optical), Sentinel-1 (radar), CHIRPS (rainfall), ERA5-Land (climate)
and NASADEM (terrain) were selected as the five sources.

\item \textbf{Google Earth Engine.} Composites were built for each
window and 15 GeoTIFF files were exported for every complete patch
(366 patches; 5,490 files).

\item \textbf{Process data and create labels.} DIP produced the RGB
chips and tensors. Patch-mean NDVI change determined the automatic
label (\texttt{candidate\_class}). Spectral and environmental values
were written into a structured evidence card. Each patch became one
JSONL record: images + evidence text + \texttt{candidate\_class}.

\item \textbf{Freeze the dataset split.} Records were split
256/55/55 (train/development/test) by stratified sampling on the
label (seed 42; 70/15/15).

\item \textbf{Run three separate experiments on the same 55 test
patches.}
\begin{itemize}
\item \textbf{A. Classification.} Four VLMs received the images and
the full evidence card (Condition~C). The run was repeated after
removing the summary line that contains the label
(\texttt{evidence\_only}). The measure is class accuracy.
Supporting analyses on this track were the closed-set probe and the
$\Delta$NDVI baselines.
\item \textbf{B. Class LoRA.} Qwen2-VL-2B was fine-tuned on 16
training patches for one epoch with the classification target, then
tested on the same 55 patches. The measure is class accuracy.
\item \textbf{C. Explanations.} The label was removed from the card
and models were asked to justify a prediction. Free-form prompts,
structured cite-only prompts and the explanation LoRA were compared.
The measures are claim correctness and class accuracy.
\end{itemize}
\end{enumerate}

The development set was used only for prompt wording, the closed-set
bias and the explanation LoRA checkpoint. No headline result was
chosen by looking at the test set.
% ============================================================
% 3. RESULTS AND DISCUSSION
% ============================================================

\section{Results and Discussion}

\subsection{Dataset Characteristics}

Table~\ref{tab:data} summarises the benchmark. The optical index changes
move together: $\Delta$NDVI correlates with $\Delta$NBR, $\Delta$NDMI
and $\Delta$EVI at 0.80, 0.77 and 0.78. Mean window rainfall fell from
69.1\,mm in 2022 to 25.1\,mm in 2023, the temperature anomaly rose from
$-1.52^\circ$C to $+0.88^\circ$C, and the soil-moisture anomaly fell
from $+0.024$ to $-0.043$. Mean VV and VH backscatter changed by $-0.47$
and $-0.48$\,dB. On average, therefore, the second dry season was drier
and slightly less green, which is why ``increase'' is a relative rank
and why many increase patches have $\Delta$ close to zero.

\begin{table}[t]
\centering
\small
\begin{tabular}{lr}
\toprule
Quantity & Value \\
\midrule
Study rectangle & 2,350.7\,km$^2$ \\
Candidate pool / final patches & 1,200 / 366 \\
Decline / stable / increase & 120 / 124 / 122 \\
Train / development / test & 256 / 55 / 55 \\
Test decline / stable / increase & 18 / 18 / 19 \\
Majority baseline (test) & 34.5\% \\
Human agreement & 8/19 (42\%) \\
Mean $\Delta$NDVI (s.d.) & $-0.016$ (0.027) \\
Patches with $|\Delta\mathrm{NDVI}| \ge 0.05$ & 34 \\
Card trend words disagree with label & 108 \\
\bottomrule
\end{tabular}
\caption{EcoExplain-VLM dataset summary. Labels are NDVI-change ranks
within this region and pair of years.}
\label{tab:data}
\end{table}

\subsection{Zero-Shot Model Performance}

Table~\ref{tab:zeroshot} reports zero-shot accuracy on the test set.
Qwen2-VL-2B is the only model above the majority baseline, at 56.4\%
(31/55; 95\% interval 43.3--68.6\%). Even this margin is not
statistically secure at $n = 55$: McNemar's test against always
predicting ``increase'' gives $p = 0.12$. The other three models do not
separate the classes. LLaVA-Interleave-0.5B answered ``decrease'' for
all 55 patches, PaliGemma answered ``decline'' for 42 of 55, and BLIP-2
repeated its own prompt, from which the parser read the first class word
(``decline''). BLIP-2's 32.7\% is therefore a parsing artefact rather
than a prediction.

\begin{table*}[t]
\centering
\small
\begin{tabular}{llrcr}
\toprule
Model or baseline & Input & Acc. & 95\% CI & Macro-F1 \\
\midrule
Qwen2-VL-2B & images + full card & \textbf{56.4\%} & 43.3--68.6 & 0.466 \\
PaliGemma-3B & images + NDVI line & 36.4\% & 24.9--49.6 & 0.271 \\
LLaVA-Interleave-0.5B & images + full card & 32.7\% & 21.8--45.9 & 0.164 \\
BLIP-2 OPT-2.7B & stitched image only & 32.7\% & 21.8--45.9 & 0.164 \\
\midrule
LLaVA-Interleave-0.5B, \texttt{evidence\_only} & images + card without label line & 36.4\% & 24.9--49.6 & 0.251 \\
LLaVA-Interleave-0.5B, closed-set probe & class-word likelihood, dev-fitted bias & 60.0\% & 46.8--71.9 & 0.479 \\
\midrule
Majority class (always ``increase'') & --- & 34.5\% & 23.4--47.7 & 0.171 \\
$\Delta$NDVI rule, $\pm 0.02$ cut-offs & card $\Delta$NDVI & 63.6\% & 50.4--75.1 & 0.606 \\
$\Delta$NDVI rule, cut-offs fitted on train & card $\Delta$NDVI & 94.5\% & 85.1--98.1 & 0.945 \\
\bottomrule
\end{tabular}
\caption{Classification on the test set ($n = 55$). The top block is
generative zero-shot Condition~C. The middle block contains supporting
runs that are not zero-shot generative results. The bottom block uses no
model.}
\label{tab:zeroshot}
\end{table*}

Two observations follow. First, the result is not explained by model
size alone: Qwen2-VL-2B and LLaVA-Interleave-0.5B received the same
input, and only the larger of the two used it. Model size cannot be
compared with PaliGemma or BLIP-2 because their inputs were shorter.
Second, the full-card condition is not a pure test of visual or spectral
understanding. Qwen2-VL-2B may be reading the summary line, the index
values, the images or a mixture of these, and it has not yet been tested
on the full test set without the summary line. For LLaVA, removing the
line made almost no difference (32.7\% to 36.4\%), because it did not
use the card in either condition.

The baselines show that the card contains enough information to solve
the task: a single threshold on its $\Delta$NDVI value reaches 94.5\%
with cut-offs fitted on the training set. This is expected, since the
labels are defined from NDVI change, and it sets a clear target that
none of the VLMs approaches. The closed-set probe lifts LLaVA to 60.0\%,
which suggests that constraining the output helps a small model even
when free generation fails, although the probe never predicts
``stable''.

\subsection{Effect of Fine-Tuning}

The class LoRA raises Qwen2-VL-2B from 56.4\% to \textbf{67.3\%} (37/55;
95\% interval 54.1--78.2\%; Table~\ref{tab:lora}). The intervals of the
two runs overlap, so with 16 training patches and $n = 55$ test patches
this is an indicative gain rather than a powered comparison. The
fine-tuned model was correct on all 18 decline and all 19 increase test
patches, but it labelled every stable patch as decline. The adapter
therefore learned two of the three classes and lost the third.

\begin{table}[t]
\centering
\small
\begin{tabular}{@{}lcc@{}}
\toprule
Qwen2-VL-2B run & Train & Test acc. \\
\midrule
Zero-shot (full card) & 0 & 56.4\% \\
Class LoRA, 1 epoch (full card) & 16 & 67.3\% \\
\midrule
\multicolumn{3}{l}{\emph{Excluded or exploratory}} \\
Full card, 2--3 epochs & 256 & 100\% \\
Card in training, images-only test & 128 & 32.7\% \\
\bottomrule
\end{tabular}
\caption{Class LoRA results for Qwen2-VL-2B. The 100\% runs copied the
label from the summary line and are not counted as results.}
\label{tab:lora}
\end{table}

Two exploratory runs support this reading of the summary line. Adapters
trained on all 256 patches with the full card reached 100\% because
they could read the answer from the card. An adapter trained with the
card on 128 patches and then tested on images alone scored 32.7\%,
below the majority baseline. Neither run is a valid estimate of
performance, and the second mixes training and test conditions, but
together they suggest that the fine-tuned model relies on the card text
rather than on the images. The explanation LoRA is a
different task and is discussed in Section~\ref{sec:expl}.

\subsection{Error Analysis}

Zero-shot Qwen2-VL-2B was correct on all 18 decline patches and on 13 of
18 stable patches, but it answered ``stable'' for all 19 increase
patches. For about half of the increase class the card's trend words
say ``stable'', and the model followed those words rather than the
summary line. Errors therefore concentrate where the label is a weak
rank: increase patches with $\Delta$ near zero are reasonably called
stable, and decline and stable patches overlap slightly after the
second-batch percentile shift.

The 108 cards whose trend words disagree with the label build a
contradiction into the input: a model that reads the descriptive text
faithfully can still miss the gold label. Together with the low human
agreement (8 of 19), this means that some ``errors'' are likely to be
label noise rather than model failure. The test set was not relabelled
after seeing model outputs.

The failures of the other models are mostly failures to follow the task.
LLaVA always produced the same answer, PaliGemma refused long prompts
and largely answered with one class, and BLIP-2 echoed its instruction.
The parser accepted the echo because it searched the whole output for a
class word; a stricter parser should read only newly generated text and
reject copied instructions.

\subsection{Explanation Quality}
\label{sec:expl}

All explanation runs used the card without the label line.
Table~\ref{tab:explain} summarises them. With the first, free-form
prompt, which asked the model to describe and justify the change in its
own words, PaliGemma answered with a class or ``unanswerable'' and no
reason, and BLIP-2 repeated the prompt, so neither produced a checkable
claim. LLaVA wrote fluent but repetitive text in which 56.7\% of the
checkable claims were correct. Qwen2-VL-2B was run on only three
patches in this mode, and one of its claims misstated the rainfall
change. Free-form output from these models is not usable as an audit
trail.

\begin{table*}[t]
\centering
\small
\begin{tabular}{lrrrrr}
\toprule
Setting and model & Class acc. & Claims / expl. & Correct claims & $\ge 1$ wrong & Self-consistent \\
\midrule
Free-form, LLaVA-Interleave-0.5B & 29.1\% & 1.09 & 56.7\% & --- & --- \\
Free-form, PaliGemma-3B & 33.3\%$^{a}$ & 0 & --- & --- & --- \\
Free-form, BLIP-2 OPT-2.7B & 32.7\%$^{b}$ & 0 & --- & --- & --- \\
\midrule
Structured, Qwen2-VL-2B & 63.6\% & 2.85 & \textbf{75.8\%} & 36.4\% & 100\% \\
Structured, LLaVA-Interleave-0.5B & 45.1\%$^{c}$ & 2.33 & \textbf{71.1\%} & 34.5\% & 90.7\% \\
\midrule
Explanation LoRA, Qwen2-VL-2B & 61.8\% & 6.84 & 50.5\% & 74.5\% & 96.4\% \\
\bottomrule
\end{tabular}
\caption{Explanation results on the test set, label removed from the
card. $^{a}$12 of 36 answered patches. $^{b}$Prompt echo; every answer
parsed as ``stable''. $^{c}$23 of 51 answered patches. The free-form
Qwen2-VL-2B run (3 patches, 43.8\% correct claims) is omitted.}
\label{tab:explain}
\end{table*}

The structured prompt gives the model a numbered list of allowed facts
that must be copied verbatim as supporting evidence, a line stating an
NDVI decision rule for the patch ($\pm 0.02$), and a request for a
one-sentence justification about NDVI only. Under this prompt, claim
correctness rises to 75.8\% for Qwen2-VL-2B and 71.1\% for
LLaVA-Interleave-0.5B. Most of the remaining incorrect claims are not
invented facts: 35 of Qwen's 38 and 35 of LLaVA's 37 flagged claims
describe a copied change smaller than the scorer's dead-band (for
example, an EVI change of $-0.006$ described as a decrease). High claim
correctness here means that the model did not invent numbers; it does
not mean that the model reasoned about the change
. Qwen2-VL-2B's class matched the stated
rule on 53 of 55 patches, and the rule alone scores 63.6\%, exactly
Qwen's accuracy in this mode. This is instruction-following, not a new
classification ability.

The explanation LoRA reached 61.8\% class accuracy (34/55; 95\% interval
48.6--73.5\%) and produced more claims per explanation (6.84), but only
50.5\% of them were correct and 74.5\% of explanations contained at
least one incorrect claim. This is worse than structured zero-shot
prompting. The bottleneck is the teacher: only 28 of 256 of its own
generations passed the filter, and only 16 were used. With this
evidence, self-distillation of a 2B model does not teach faithful
explanation; the adapter mainly learned the output format.

\subsection{Comparison with Existing Approaches}

Pixel and segment-based change detectors optimize map accuracy
rather than a three-class patch rank with a sentence, and they remain
the right tool for a wall-to-wall disturbance map. Remote-sensing
question answering  and change captioning
 pair imagery with language, but they work on
single-sensor optical imagery and do not check generated text against
measured values. EcoExplain-VLM asks a different question: whether a
general VLM can use the same multi-sensor evidence that a human analyst
would see, and whether its explanation agrees with that evidence.

Remote-sensing VLMs such as GeoChat are trained on
overhead imagery and spatial referring expressions and are a natural
next baseline. They were not evaluated here because the available
compute limited the study to general checkpoints of 0.5--3 billion
parameters. CLIP-style encoders could score the
image pair without generating text, but they would not produce the
justification that the benchmark is designed to test.

The most informative comparison this dataset supports is with simple
baselines. Against the majority class, only Qwen2-VL-2B with the full
card shows a clear margin. Against a single $\Delta$NDVI threshold, none
of the models comes close.

% ============================================================
% 4. CONCLUSIONS
% ============================================================

\section{Conclusions}

EcoExplain-VLM packages 366 Kanha patches, matched 2022/2023 dry-season
Sentinel-2 and Sentinel-1 data, climate and terrain context, and a
three-class NDVI-change label into a reproducible VLM benchmark. On the
55-patch test set, Qwen2-VL-2B reaches 56.4\% zero-shot accuracy with
the full evidence card, and a 16-patch class LoRA raises this to
67.3\%; both figures depend on a card that contains the candidate label.
LLaVA-Interleave-0.5B, PaliGemma-3B and BLIP-2 OPT-2.7B stay at
32.7--36.4\%, the level of always guessing one class, while a single
$\Delta$NDVI threshold reaches 94.5\%. For explanations, free-form
prompts give unusable output, and a structured cite-only prompt makes
explanations consistent with the evidence (75.8\% and 71.1\% correct
claims) at the cost of restating the supplied facts. An explanation LoRA
trained on 16 self-generated justifications gives 61.8\% class accuracy
and 50.5\% correct claims and is a pilot rather than an improvement.

Small open VLMs, as evaluated here, are not a substitute for spectral
change detection in Kanha. Under a constrained prompt they can repeat
the facts they are given. The value of the benchmark is that it makes
this limitation measurable and gives larger or domain-adapted models a
fixed, label-free test against which to show real progress.

% ============================================================
% LIMITATIONS (required by ACL; unnumbered, after the conclusion)
% ============================================================

\section*{Limitations}

\paragraph{Labels.} The classes are percentile ranks of one-year
dry-season NDVI change, not validated disturbance, degradation or
recovery, and the absolute changes are small. Human agreement on 19
patches is 42\%, with a single annotator. An expert audit with a second
annotator is needed before any management use.

\paragraph{Label in the input.} The Condition~C card contains
\texttt{candidate\_class}, and both headline Qwen2-VL-2B results (56.4\%
and 67.3\%) use it. Future evaluations should use the
\texttt{evidence\_only} card by default, which the code now supports.


\paragraph{Sample size and split.} With $n = 55$, 95\% intervals are about
$\pm 13$ percentage points wide. The random stratified split does not
prevent spatial leakage ; a spatial block split should
be added as a second test set.

\paragraph{Models.} The results cover four small open VLMs and say nothing
about larger proprietary or remote-sensing-pretrained models.

\paragraph{Explanation scoring.} Automatic claim matching can miss
paraphrases, can pair a direction word with the wrong indicator, and can
mark a true but small change as incorrect because of the dead-band. The
reported claim correctness is not a human judgement of faithfulness.



\paragraph{Future directions.} Five extensions would change the scientific picture rather than add
rows to a leaderboard: (i)~make label-free evaluation the only primary
table, including Qwen2-VL-2B on all 55 test patches; (ii)~add a spatial
block test; (iii)~collect expert labels for all 366 patches or a
sufficiently large subsample; (iv)~fine-tune on all 256 training
patches without the label line, with early stopping on the development
set; and (v)~evaluate a larger or remote-sensing-pretrained VLM, and a
stronger teacher for explanation distillation, on the same cards.



% ============================================================
% REFERENCES
% ============================================================
\begin{thebibliography}{10}

\bibitem{wang2024qwen2vl}
Peng Wang et~al. 2024.
\newblock Qwen2-VL: Enhancing vision-language model's perception of the world at any resolution.
\newblock \emph{arXiv preprint arXiv:2409.12191}.

\bibitem{li2024llavainterleave}
Feng Li et~al. 2024.
\newblock LLaVA-NeXT-Interleave: Tackling multi-image, video, and 3D in large multimodal models.
\newblock \emph{arXiv preprint arXiv:2407.07895}.

\bibitem{beyer2024paligemma}
Lucas Beyer et~al. 2024.
\newblock PaliGemma: A versatile 3B VLM for transfer.
\newblock \emph{arXiv preprint arXiv:2407.07726}.

\bibitem{li2023blip2}
Junnan Li, Dongxu Li, Silvio Savarese, and Steven Hoi. 2023.
\newblock BLIP-2: Bootstrapping language-image pre-training with frozen image encoders and large language models.
\newblock In \emph{Proceedings of ICML}.

\bibitem{hu2022lora}
Edward~J. Hu et~al. 2022.
\newblock LoRA: Low-rank adaptation of large language models.
\newblock In \emph{Proceedings of ICLR}.

\bibitem{lobry2020rsvqa}
Sylvain Lobry, Diego Marcos, Jesse Murray, and Devis Tuia. 2020.
\newblock RSVQA: Visual question answering for remote sensing data.
\newblock \emph{IEEE Transactions on Geoscience and Remote Sensing}, 58(12):8555--8566.

\bibitem{liu2022changecaption}
Chenyang Liu, Rui Zhao, Hao Chen, Zhengxia Zou, and Zhenwei Shi. 2022.
\newblock Remote sensing image change captioning with dual-branch transformers: A new method and a large scale dataset.
\newblock \emph{IEEE Transactions on Geoscience and Remote Sensing}, 60:1--20.

\bibitem{li2023pope}
Yifan Li, Yifan Du, Kun Zhou, Jinpeng Wang, Wayne~Xin Zhao, and Ji-Rong Wen. 2023.
\newblock Evaluating object hallucination in large vision-language models.
\newblock In \emph{Proceedings of EMNLP}.

\bibitem{jacovi2020faithfulness}
Alon Jacovi and Yoav Goldberg. 2020.
\newblock Towards faithfully interpretable NLP systems: How should we define and evaluate faithfulness?
\newblock In \emph{Proceedings of ACL}, pages 4198--4205.

\bibitem{turpin2023unfaithful}
Miles Turpin, Julian Michael, Ethan Perez, and Samuel~R. Bowman. 2023.
\newblock Language models don't always say what they think: Unfaithful explanations in chain-of-thought prompting.
\newblock In \emph{Advances in NeurIPS}.

\end{thebibliography}

% ============================================================
% APPENDIX
% ============================================================

\clearpage
\appendix

\section{Additional Experimental Details}

\subsection{Fine-Tuning Configuration}

Table~\ref{tab:ftconfig} lists the settings shared by the class and
explanation LoRA runs. Both used Qwen/Qwen2-VL-2B-Instruct with the
Hugging Face \texttt{transformers} and \texttt{peft} libraries, and both
trained for one epoch on 16 examples. For the class LoRA, the 16
patches were the configured training size. For the explanation LoRA,
the 16 were the filtered justifications that fell within the first 128
training patches.

\begin{table}[ht]
\centering
\small
\begin{tabular}{@{}ll@{}}
\toprule
Setting & Value \\
\midrule
Base model & Qwen2-VL-2B-Instruct \\
Rank / $\alpha$ / dropout & 8 / 16 / 0.1 \\
Target modules & q, k, v, o, gate, up, down \\
Optimiser & AdamW \\
LR / weight decay & $5\times10^{-5}$ / 0.05 \\
Gradient accumulation & 4 steps \\
Epochs / examples & 1 / 16 \\
Hardware & GTX 1080 Ti (Ada cluster) \\
\bottomrule
\end{tabular}
\caption{LoRA configuration.}
\label{tab:ftconfig}
\end{table}

\subsection{Prompt and Output Parsing}

\paragraph{Classification.} Qwen2-VL-2B and LLaVA-Interleave-0.5B
received the two images, the full evidence card and the instruction
``Reply with exactly one word after the images: decline or increase or
stable. Do not copy the evidence text.'' PaliGemma received the two
images, the card's NDVI line and ``Two forest RGB chips: first image T0,
second T1. Vegetation change class --- answer with exactly one word:
stable, decline, or increase.'' BLIP-2 received the stitched image and
``Side by side: left forest patch T0, right forest patch T1. One word
only --- decline, increase, or stable:''. The \texttt{evidence\_only}
variant removes every card line that names \texttt{candidate\_class}.

\paragraph{Explanations.} The free-form prompt tells the model that
image 1 shows the patch in January--February 2022 and image 2 in
January--February 2023, asks it to describe the visible change in one
sentence, decide whether vegetation is stable, declined or increased,
and justify the decision. Qwen2-VL-2B and LLaVA answer in JSON;
PaliGemma and BLIP-2 receive a one-sentence version of the request.
The structured prompt adds a numbered list of allowed facts to be copied verbatim, a line stating
the $\pm 0.02$ NDVI rule for the patch, and a one-sentence justification
limited to NDVI. An automatic check confirms that no explanation prompt
contains \texttt{candidate\_class}.

\paragraph{Parsing.} The parser lowercases the output, accepts a
single-word answer directly, and otherwise searches for the first class
word or synonym. It does not exclude text copied from the prompt, which
explains BLIP-2's echo artefact.

\paragraph{Teacher filtering.} For the explanation LoRA, a generation
was kept only if its parsed class equalled the label, it contained at
least two checkable claims, and every claim agreed with the measured
change. This filter kept 28 of 256 training patches.

\subsection{Additional Results}

With $n = 55$, intervals are wide: 31/55 (56.4\%) corresponds to
43.3--68.6\% and 37/55 (67.3\%) to 54.1--78.2\%. Per-class counts
(about 18 patches per class) are too small to support a detailed error
taxonomy. The closed-set LLaVA probe (60.0\%) relies on a
development-fitted bias and is reported as a probe, not as a zero-shot
result; the same probe left PaliGemma and BLIP-2 at 32.7\%, predicting
``stable'' for every patch. In the \texttt{evidence\_only} runs, LLaVA
answered ``decline'' for 51 patches and ``increase'' for 4.

\end{document}
